\documentclass[11pt]{amsart}
\usepackage{amssymb,booktabs,hyperref}
\newtheorem{theorem}{Theorem}
\newtheorem{lemma}[theorem]{Lemma}
\newcommand{\MS}{\mathrm{MS}}
\newcommand{\AK}{\mathrm{AK}}

\title{Andrews--Curtis certificates for Miller--Schupp presentations}
\author{peteracs}
\date{November 2026}

\begin{abstract}
Of the $1190$ Miller--Schupp presentations $\MS(n,w)$ with $n\le 7$ and
$|w|\le 7$, a published greedy search left $657$ unsolved. We give explicit,
machine-checkable sequences of elementary Andrews--Curtis moves that
trivialize $78$ of them, and that show $65$ more to be AC-trivial if and only
if the Akbulut--Kirby presentation $\AK(n)$ is; for $n=3$ this adds $27$
presentations to those known to stand or fall with $\AK(3)$. A further $17$
are tied to two other presentations with $n=7$. Every certificate is checked
by a program of about 150 lines that uses only free reduction.
\end{abstract}

\maketitle

\section{Setting}

An \emph{AC move} on a balanced presentation $\langle x,y\mid r_0,r_1\rangle$
replaces $r_i$ by $r_ir_j^{\pm1}$ ($j\ne i$), by $r_i^{-1}$, or by
$gr_ig^{-1}$ for a letter $g$, and then freely reduces. The Andrews--Curtis
conjecture asserts that every balanced presentation of the trivial group is
AC-equivalent to $\langle x,y\mid x,y\rangle$. The Miller--Schupp
presentations
\[
\MS(n,w)=\langle x,y\mid x^{-1}y^nx=y^{n+1},\ x=w\rangle,
\]
with $w$ of exponent sum zero in $x$, present the trivial group and are a
standard source of potential counterexamples, as are the Akbulut--Kirby
presentations $\AK(n)=\langle x,y\mid x^n=y^{n+1},\ xyx=yxy\rangle$, open for
$n\ge3$. Shehper et al.~\cite{shehper} studied the $1190$ presentations with
$n\le7$, $|w|\le7$; their greedy search solved $533$. They proved
(Theorem~B) that $\MS(1,w)$, $\MS(n,y^{-1}xyx^{-1})$ and $\MS(2,w_k)$ with
$w_k=y^{-k}x^{-1}yxy$ are AC-trivial, that $\MS(n,w_k)$ is AC-equivalent to
$\AK(n)$, and (Theorem~10) trivialized two more. Lisitsa~\cite{lisitsa}
trivialized $\MS(n,y^{-1}xyx^{-1})$ for $3\le n\le8$.

\section{Method}

The search works on pairs of cyclic words, each up to rotation and
inversion, with two steps: replace one relator by the cyclic reduction of
$uv$, with $u$ a rotation of it and $v$ a rotation of the other relator or of
its inverse; or apply one of the Whitehead automorphisms $y\mapsto yx^{\pm1}$,
$x\mapsto xy^{\pm1}$ to both relators when this shortens the pair. Products
are compositions of AC moves. Automorphisms are not, and are removed as
follows.

\begin{lemma}
Let $P$ be a presentation, $\varphi\in\mathrm{Aut}(F(x,y))$, and suppose
$\varphi(P)$ is carried to $Q$ by AC moves $m_1,\dots,m_k$. Then $P$ is
carried to $\varphi^{-1}(Q)$ by the AC moves $\varphi^{-1}(m_1),\dots,
\varphi^{-1}(m_k)$, where $\varphi^{-1}$ fixes multiplication and inversion
moves and sends conjugation by $g$ to conjugation by the word
$\varphi^{-1}(g)$, a composition of conjugations by letters.
\end{lemma}

Applying the lemma at each automorphism, a search path from $P$ to
$\langle x,y\mid x,y\rangle$ becomes AC moves from $P$ to
$\langle x,y\mid\psi(x),\psi(y)\rangle$ for an automorphism $\psi$; since
$\{\psi(x),\psi(y)\}$ is a basis, Nielsen moves (themselves AC moves) take it
to $\langle x,y\mid x,y\rangle$. A path from $P$ to a presentation $K$ becomes
AC moves from $P$ to $\psi(K)$, and $\psi(K)$ is AC-trivial if and only if $K$
is, by the lemma again. Ties to a class are found with a saved search graph:
the graph is grown from $K$, and a search from $P$ stops when it reaches a
state of the graph.

\section{Results}

\begin{center}
\begin{tabular}{rrrrrr}
\toprule
$n$ & AC-trivial & tied to $\AK(n)$ & tied to other & already known & open\\
\midrule
2 & 0 & 0 & 0 & 0 & 32\\
3 & 1 & 27 & 0 & 8 & 36\\
4 & 23 & 15 & 0 & 8 & 72\\
5 & 31 & 9 & 0 & 7 & 88\\
6 & 15 & 7 & 0 & 7 & 117\\
7 & 8 & 7 & 17 & 7 & 115\\
\midrule
all & 78 & 65 & 17 & 37 & 460\\
\bottomrule
\end{tabular}
\end{center}

\noindent The columns count the $657$ presentations the greedy search left.
``Already known'' are those covered by the results recalled above. The other
two cores are $\MS(7,yxy^{-1}x^{-1})$ (13 tied) and $\MS(7,yxy^2x^{-1})$
(4 tied); both are open. The trivializations use 122 to 10868 moves (median
576). No presentation with $n=2$ was resolved, and $\AK(3)$ was not
trivialized.

The comparison with~\cite{shehper} is exact for the families and theorems
above. That paper also relates four presentations of length 14 to $\AK(3)$
without listing them in machine-readable form; one tie here has length 14
and may coincide with one of them.

\section{Data}

The certificates, the table and the checker are at
\url{https://github.com/peteracs/ac-miller-schupp}. A certificate lists the
moves on the ordered pair of relators; the checker replays them from the
stated start and compares with the stated end.

\begin{thebibliography}{9}
\bibitem{shehper} A.~Shehper et al., \emph{What makes math problems hard for
reinforcement learning: a case study}, arXiv:2408.15332v2, 2025.
\bibitem{lisitsa} A.~Lisitsa, \emph{New Andrews--Curtis trivializations for
Miller--Schupp group presentations}, Examples and Counterexamples 6 (2024),
100168.
\end{thebibliography}

\end{document}
